\chapter{Conclusion}
\label{chap:conclusion}

% ----------------------------------------------------------------
\section{Bilan de la recherche}

Ce mémoire s'est donné pour objectif d'identifier, de mesurer et de réduire les
goulots de performance de MetaDONE, un outil MetaCASE bâti sur l'architecture
MetaL$_2$ et reposant sur une couche de persistance JPA. La démarche adoptée a
suivi un cycle empirique complet pour chacun des quatre scénarios étudiés~:
conception d'un benchmark instrumenté, mesure de référence, identification de
la cause racine dans le code source, conception et implémentation d'une
optimisation, vérification expérimentale du gain.

Cette méthodologie a guidé toutes les décisions techniques, y compris les
décisions négatives. Sur les quatre optimisations candidates initialement
identifiées, deux ont été \emph{rejetées} sur la base des mesures~:
la candidate \og batch-JPA \fg{} (activation du paramètre \texttt{batchLimit}),
sans effet sur une base H2 embarquée, et la candidate \og bypass-précondition \fg{}
(introduction d'une méthode \texttt{createRoleTrusted} contournant
\texttt{precondition\_createRole}), incompatible avec le principe Design by
Contract structurant l'architecture de MetaDONE. %Ces refus documentés délimitent l'espace des solutions acceptables et constituent à part entière un résultat de la recherche.

% ----------------------------------------------------------------
\section{Contribution}

La contribution technique principale de ce travail se résume en deux
optimisations algorithmiques implémentées et validées.

L'optimisation du Scénario~1 introduit un cache \texttt{WeakHashMap} dans
\texttt{ConcreteModel\_impl.\_isInDefinition}, éliminant un scan linéaire
répété de la liste de définition. La complexité de \texttt{getMOByName} passe
de $O(N^2)$ à $O(N)$.

L'optimisation Opt-5A du Scénario~5 corrige une asymétrie non documentée entre
\texttt{\_createObject} (insertion directe en $O(1)$) et \texttt{\_createRole}
(passage par un scan en $O(\text{taille\_du\_modèle})$). En quatre lignes
modifiées, la création d'arcs passe de $O(N^2)$ à un plancher proche de
$O(N \log N)$ pour les topologies à degré borné, avec un gain médian de
$2{,}19\times$ à $N = 1\,000$ et $4{,}39\times$ à $N = 5\,000$ (mesuré sur
10 répétitions par campagne). La borne résiduelle correspond au plancher inhérent à tout
backend de persistance indexé.

Au-delà des chiffres, ce travail dégage trois enseignements généralisables.
Premièrement, les bottlenecks les plus impactants ne se trouvent pas
toujours là où l'intuition les place : la mise en évidence de la complexité
$O(N^2)$ de \texttt{\_createRole} n'a été possible qu'en comparant deux
topologies dont l'une infirmait l'hypothèse initiale. Deuxièmement, dans
un système à garanties formelles, les optimisations les plus sûres sont celles
qui \emph{exploitent} les invariants plutôt que celles qui les contournent~:
l'opposition entre la candidate \og bypass-précondition \fg{} (rejetée) et
Opt-5A (retenue) en est l'illustration.
Troisièmement, l'explicitation des hypothèses simplificatrices — ici,
l'exclusion délibérée de l'optimisation du schéma JPA — est essentielle à
l'interprétation correcte des résultats.

% ----------------------------------------------------------------
\section{Limites du travail et points de vigilance}

Plusieurs limites doivent être signalées pour délimiter la portée des résultats présentés.

\paragraph{Périmètre instrumenté.}
Les benchmarks ont été réalisés en mode standalone JUnit, hors du contexte
Eclipse RCP. Si ce choix offre une excellente contrôlabilité pour les couches
métier et repository (Scénarios 1, 2 et 5), il a empêché de mesurer
directement l'étape~3 du pipeline de visualisation (évaluation des expressions
GraSyLa). Cette limite a été compensée a posteriori par un profilage live de
l'application Eclipse via async-profiler (Section~\ref{sec:sc4-live}), mais
les deux types de mesures ne sont pas strictement comparables et la
quantification reste indirecte.

\paragraph{Biais d'accumulation JPA.}
La base de test étant partagée entre exécutions, la session OpenJPA accumule
des entités au fil des scénarios. Ce phénomène (\emph{entity manager bloat}) a
été observé à $N = 5\,000$ en S5 et a légèrement biaisé les valeurs absolues
mesurées. Les comparaisons relatives avant/après optimisation restent valides,
mais les valeurs absolues à grand $N$ doivent être interprétées avec
prudence. Une réexécution sur base fraîche pour chaque scénario améliorerait
la fiabilité des chiffres absolus.

\paragraph{Hypothèse simplificatrice sur la persistance.}
L'optimisation du schéma JPA et des index SQL a été explicitement exclue du
périmètre. Cette exclusion, motivée par la portée architecturale d'une telle
modification, signifie que le plancher proche de $O(N \log N)$ observé en S5
après Opt-5A reflète une contrainte du backend (mise à jour d'index B-tree à
chaque INSERT) non remise en cause dans ce travail.

\paragraph{Si c'était à refaire.}
Avec le recul, deux choix méthodologiques mériteraient d'être révisés. D'une
part, le profilage live aurait dû être réalisé plus tôt dans le processus :
les 11~\% de coût attribués à \texttt{getEquations}/\texttt{hasType} en
contexte interactif (contre $\sim 0\,\%$ en standalone) n'ont été révélés
qu'en fin de travail et auraient pu orienter différemment les priorités.
D'autre part, l'effort consacré à l'évaluation des candidates batch-JPA et
bypass-précondition (finalement rejetées) aurait pu être mieux dimensionné
avec une lecture préalable plus attentive de la configuration H2 et du
contrat DbC.
Ces deux constats relèvent toutefois davantage d'un apprentissage que d'une
remise en cause de la démarche d'ensemble.

% ----------------------------------------------------------------
\section{Pistes non explorées et perspectives}

Plusieurs pistes de réflexion identifiées au cours du travail n'ont pu être
développées, soit par manque de temps, soit par dépassement du périmètre
défini. Elles méritent toutefois d'être signalées pour orienter des travaux
ultérieurs.

\paragraph{Mise en cache de \texttt{getEquations()} — court terme.}
Le profilage live a mesuré que cette méthode représente 11,16~\% du temps CPU
en contexte interactif, car elle reconstruit l'intégralité des 106 wrappers
\texttt{GrasylaEquation} à chaque cycle de rendu. Un mécanisme de cache
volatile invalidé par \texttt{notifyScriptChanged()} suffirait à éliminer ce
surcoût en quelques lignes de code. C'est la piste au rapport gain/effort le
plus favorable parmi celles identifiées.

\paragraph{Index de précondition pour \texttt{MANY\_TO\_MANY} — moyen terme.}
Le bottleneck résiduel identifié en S5 pour les topologies à fort degré
(complexité $O(N^2)$ due à la vérification de l'unicité des paires) pourrait
être corrigé par un index $\texttt{Map<MetaRole, Map<K\_Object,
Set<K\_Object>>>}$ maintenu dans \texttt{ConcreteModel\_impl}. L'implémentation
requiert une gestion de cohérence à la suppression et une initialisation
\emph{lazy} pour les modèles chargés.

\paragraph{Optimisation de la couche de persistance — long terme.}
L'introduction d'index composites SQL sur \texttt{(Property\_domain,
Property\_type)} convertirait certains scans $O(\text{degré})$ en $O(\log N)$.
Cette piste relève cependant d'un projet à part entière, interagissant avec
l'ensemble des backends supportés. Le promoteur a souligné que, dans le scope
de ce mémoire, l'option ORM/SQL restant la voie de persistance par défaut,
l'optimisation de la base de données constituera \emph{une étape inévitable}
pour tout travail ultérieur visant un passage à l'échelle significatif.

\paragraph{Reconstruction incrémentale de la scène GraSyLa.}
Le profilage live a révélé que les $\approx 53\,\%$ de temps CPU passés dans la
boucle \texttt{trigger}~$\rightarrow$~\texttt{buildComponent}~$\rightarrow$~%
\texttt{evaluateComponent} proviennent d'une reconstruction \emph{globale} de
la scène à chaque modification du modèle. Une stratégie de reconstruction
incrémentale, ciblée sur les seuls composants affectés, semble la piste à plus
fort potentiel pour les améliorations futures de l'interface graphique.

% ----------------------------------------------------------------
\section*{}

MetaDONE est un outil dont l'architecture privilégie l'extensibilité, la
généricité et la correction formelle ; cette généricité a un coût, qu'aucune
optimisation locale ne peut entièrement effacer. Le travail présenté ici ne
prétend pas refonder cette architecture, mais en améliorer le comportement
empirique aux points où une analyse rigoureuse révèle une marge d'amélioration
réelle. Les deux optimisations retenues — toutes deux locales, peu invasives,
et préservant intégralement les garanties du Design by Contract — démontrent
qu'un gain substantiel de performance est compatible avec le respect des
principes architecturaux fondateurs de l'outil.
